
The contents of this chapter are published in~\citet{visani_h-packer_2024}, under the title
\textit{H-Packer: Holographic Rotationally Equivariant Convolutional Neural Network for Protein
Side-Chain Packing}.

Determining amino acid side-chain conformations in a protein, known as Protein Side-Chain Packing
(or Rotamer Packing), is an essential step in protein folding and the de-novo design of proteins.
Computational approaches to protein folding often divide the structure inference problem into two
steps: first, they characterize the rigid backbone structure, and then they pack the side-chains
associated with the amino acids at each residue. The flexibility of the side-chain makes the search
in the space of possible conformations inevitably complex and computationally expensive. The
de-novo protein design protocols also rely on similar logical steps: often an amino acid sequence
compatible with a desirable backbone structure is to be inferred
(designed)~\citep{dauparas_robust_2022} and then the associated side-chains should be packed to
form the full atomic composition of a protein. Furthermore, it has been suggested that models
trained to learn \textit{distributions} of side-chain conformations can be used to predict
mutational effects in a zero-shot fashion~\citep{luo_rotamer_2023}, thereby making rotamer packing
a valuable proxy problem for resolving structure-to-function maps.

Many of the conventional methods for side-chain packing rely on physical models through which they
find a rotamer that minimizes a physically-reasoned heuristic energy of the protein
fold~\citep{alford_rosetta_2017, huang_faspr_2020, krivov_improved_2009,
cao_improved_2011, liang_fast_2011}. Popular algorithms include RosettaPacker from the rosetta
suite~\citep{alford_rosetta_2017}, FASPR~\citep{huang_faspr_2020}, and
SCWRL~\citep{krivov_improved_2009}.
However, these computational methods often lack accuracy and speed in their predictions.
As deep learning makes
strides in protein science, there is a growing effort in developing machine learning methods for
rotamer packing. Among these methods is DLPacker~\citep{misiura_dlpacker_2022}, which treats the
packing problem as an image transformation. This algorithm characterizes the local environment of a
given amino acid backbone within a structure as a 3D image, and uses this model to predict the
atomic coordinates of the side chain. It then compares the predicted side-chain to a pre-set
library of rotamers to select the closest conformation. AttnPacker~\citep{mcpartlon_end--end_2023},
a more recently developed method, uses a deep graph attention network to model the local geometry
of a residue within a structure and is trained to predict the coordinates of the side-chain atoms.
Recently, diffusion models over side-chain torsional angles have also being applied, such as
DiffPack~\citep{zhang_diffpack_2023}.

The specific conformation - or the ensemble of conformations - adopted by a side-chain is fundamentally
determined by the physical forces acting upon it, among which those coming from local sources are the
most prominent. Therefore, modeling a side-chain's conformation as a function of its local atomic environment
is a physics-appropriate approximation. Furthermore, the conformation can be described unambiguously by 
SO(3)-invariant or SO(3)-equivariant variables. Thus, side-chain packing fits nicely within our framework of
rotationally equivariant modeling of residue-centered protein neighborhoods.

Here, we tackle the problem of side-chain packing by learning to directly regress over $\chi$
(torsional) angles, which are main degrees of freedom determining side-chain conformations. We
derive and discuss three possible parameterizations of the $\chi$ angles, ultimately settling on
regressing over the Sine and Cosine transforms of the angles, which are rotationally invariant. We
introduce Holographic-Packer (H-Packer), a deep learning method that packs rotamers by first
predicting candidate $\chi$ angles from backbone and sequence, and then refines the predictions
with a model trained on full-atom structures.

By directly predicting the side-chain $\chi$-angles, H-packer does not rely on comparing its output
with a pre-set library of rotamers, making it computationally more efficient than methods like
DLPacker. Furthermore, H-Packer is light-weight (2$\times$3M parameters vs. 208M of AttnPacker) and
requires few resources to train (single vs. multiple GPUs for diffusion models like DiffPack). We
evaluate the packing performance of H-Packer on standard datasets, and show that it has generally
better performance than conventional physics-based methods, and competitive against machine
learning solutions. In general, our results suggest that H-Packer has learned complementary
features to alternative methods.
Code for H-Packer is freely available at \texttt{https://github.com/gvisani/hpacker}.



\section{Methods}

\subsection{Modeling side-chain conformations with $\chi$ Angles}

While amino acids are composed of a maximum of 10 heavy atoms (in the case of Tryptophan), their 3D
conformations can be uniquely described by the value of at most 4 dihedral angles, referred to as
the $\chi$ angles (Figure~\ref{fig:hpacker}A). This reduction in the number of degrees of freedom
is granted due to the physical constraints posed on the remaining internal coordinates (bond
angles, bond lengths, and dihedral angles - \textit{redundant internal coordinates}). Specifically,
the inter-atomic physical interactions within amino acids often constrain these redundant
coordinates to a constant, or a well-defined function of the residue's $\chi$ angles. Therefore,
predicting $\chi$ angles is the key step in side-chain packing.

H-packer addresses the side-chain packing problem in two steps: (i) it predicts the value of $\chi$
angles, and (ii) it reconstructs the atomic coordinates using the predicted $\chi$ angles and the
constrained values of the redundant internal coordinates. Specifically, we evaluate the redundant
internal coordinates from a subset of training data (1,700 structures) by leveraging the
\texttt{internal\_coords} feature of the \texttt{biopython} package. We empirically verified that
the distributions of values were Gaussian with low variance, and resolved to take their
\textit{medians} as the ground truth. Substituting these values for the original ones yields a
negligible Null Reconstruction error of approximately
0.127\AA\,(Figure~\ref{fig:null_reconstruction}). Notably, this error remains unchanged even when
using only 100 reference structures instead of 1,700
(Figure~\ref{fig:null_reconstruction_comparison}).

\subsection{Predicting $\chi$ angles using H-Packer}
We aim to predict the $\chi$ angles associated with a side-chain conformation from the
configuration of atoms surrounding a given residue. This atomic neighborhood is associated with the
backbone and the side-chain of the neighboring residues in the structure.

During inference only the coordinates of the backbone atoms are known a priori - alongside the
identity of the amino acids they belong to. However, physical interactions with the atoms of
neighboring side-chains are the true determinants of a residue's conformation. Therefore, we
develop H-Packer into a two-step solution (Figure~\ref{fig:hpacker}C). Specifically, we two train
models: one to predict $\chi$ angles from the backbone atoms and amino-acid identity alone, the
another to predict $\chi$ angles from full neighborhoods, i.e., by including the true side-chain
atoms of the surrounding residues (minus the residue of interest). At inference time, we use the
first model to make an \textit{initial guess} of the side-chain conformations, and then a second
model to iteratively \textit{refine} the predictions.

To build the individual models that predict $\chi$ angles, we start by considering their
symmetries. Notably, $\chi$ angles are \textit{invariant} to rigid-body transformations
(translations and rotations) of the protein (i.e., they are SE(3) invariant). Translation
invariance can be satisfied by choosing a well-defined center for a residue of interest (e.g. the
residue's C$\alpha$), leaving rotational (i.e. SO(3) invariance) about the chosen center to be
accounted for. The rotational symmetry and the dependence on local interaction make this problem
fall neatly within our framework of SO(3)-equivariant modeling of protein neighborhoods. Thus we
build two SO(3)-equivariant models - using the components described in Chapter~\ref{sec:background}
- to predict a residue's $\chi$ angles, considering as input the point cloud of atoms within a
radius $r=10\AA$ of the residue's C-$\alpha$ (with or without the neighboring side-chains).


\begin{figure}[t]
    \centering
    \resizebox{1.0\textwidth}{!}{
    \includegraphics{Quals/figures/hpacker.pdf}
    }
    \figcaption{Overview of H-Packer.}{}
    \medskip
    \small
    \justifying
    \textbf{A:} Illustration of Glutamine's $\chi$ angles, of which there are three. \textbf{B:}
    Schematic shows the SO(3)-Equivariant network for side-chain packing by first predicting the
    missing residue's $\chi$ angles from its surrounding atomic environment, and then using the
    $\chi$ angles to reconstruct the residue's side-chain. As illustrated in \textbf{C}, H-Packer
    consists of two networks, one trained on backbone atoms only and used to make an initial guess,
    and one trained on full side-chain neighborhoods and used to refine the predictions.
    \label{fig:hpacker}
    % \vspace{-4mm}
\end{figure}


\textbf{Holographic Encoding.}
We encode neighborhoods into holograms as described in Sec.~\ref{sec:encoding}. We incorporate
atom-level input features by dividing the holographic encoding into different \textit{channels}
(see Figure~\ref{fig:hpacker}). We consider the following two sets: (i) \textbf{Atomic channels}:
C, N, O, S, wildcard element excluding hydrogens, partial charge from the Amber99sb force
field~\citep{ponder_force_2003}, and (ii)  \textbf{Amino-Acid channels}: one for each of the 20
canonical amino-acids, plus a wildcard channel. We include the charge value in its dedicated
channel as the weights $\omega_{i}$ coupled to the point cloud's density function.
While we train the \textit{initial guess} model using both sets of channels (atomic and amino-acid)
as input, we only consider the atomic channels for the \textit{refinement} model. We do this in an
effort to make the model's predictions more grounded in physical interactions. We condition both
models with the identity of the residue of interest by concatenating a linear embedding of its
one-hot encoding to the input's invariant ($\ell=0$) features. This is particularly necessary for
the refinement model - which is trained only with atomic channels - since it wouldn't otherwise
know about the identity of the residue of interest.

\textbf{SO(3)-Equivariant neural network architecture.}
We use the resulting holograms as inputs to an SO(3)-Equivariant Convolutional Neural Network
(Figure~\ref{fig:hpacker}B). Our model is conceptually divided into three parts:\\
\textbf{First,} a linear layer that projects data and conditioning to a hidden representation with
same number of features per $\ell$.\\
\textbf{Second,} a stack of equivariant blocks connected via additive skip connections, each
composed of: (i) feature-wise tensor product nonlinearity, (ii) layer norm with silu nonlinearity,
and (iii) a linear layer whose output dimensions are the same as the input's (see
Sec.~\ref{sec:layers}). After the final block, we retain only the features of type $\ell = 0$ or
$\ell = 1$ depending on the training objective (Section~\ref{sec:training_objective}). It should be
noted that features of type $\ell=0$ are rotationally invariant scalars, whereas those associated
with $\ell=1$ are equivariant vectors that transform consistently with the input under rotation. We
use $\ell=1$ features to directly learn the orientation of the intersecting planes that define a
side-chain's dihedral angles $\chi$ (see Section~\ref{sec:training_objective}).\\
\textbf{Third,} optionally and only for the models with invariant ($\ell = 0$ output), we apply a
standard feed-forward neural network with dropout regularization and SiLU nonlinearity.



\subsection{Training objectives to infer $\chi$ angles}
\label{sec:training_objective}

We consider three alternative parameterizations of $\chi$ angles, i.e. three possible objective
functions:

{\bf  (i) The angle itself.} $\chi$ angles are defined between $-180\degree$ and $180\degree$ with
a periodicity such that the angles $-179\degree$ and $179\degree$ are to be considered $2\degree$
apart, not $358\degree$. Thus, plain MSE loss would pose strong and unnatural constraints on the
model. To account for this, we mod the predictions to fall in the valid range, and compute the loss
between two angles as the minimum between the computed error and $360\degree$ minus the error,
resulting in the following loss function:
\EQ
    \mathcal{L}_{\text{angles}}(\{\hat{\chi_i}\}_{i=1}^{N_{\chi}}, \{\chi_i\}_{i=1}^{N_{\chi}}) =
    \frac{1}{N_\chi} \sum_{i=1}^{N_{\chi}} \min(E_{\chi_i}, \, 2\pi - E_{\chi_i}) \quad
    \text{where} \quad E_{\chi_i} = (\text{mod}(\hat{\chi_i}, 2\pi) - \chi_i)^2
\EE
where $\hat{\chi_i}$ and $\chi_i$ are  the predicted and the true values of the $i^{th}$ $\chi$,
respectively, and $N_\chi$ is the number of $\chi$ angles associated with the residue of interest.
In our implementation, the $\chi$ angle domain is scaled and shifted to fall in $[0, 2]$ to make
the scale of the loss functions comparable between the three representations of the angles.

{\bf (ii) Sine and Cosine transforms of the angle.} A pair of sine and cosine transformation
provides an alternative representation for a $\chi$ angle that accounts for its periodicity and is
also rotationally invariant; a similar approach is also considered in concurrent
work~\citep{mukhopadhyay_zymepacknet_2023}. We directly predict sine and cosine values by feeding 8
outputs from the network to a tanh activation function, which then form the arguments of  a MSE
loss function:
\EQ
\mathcal{L}_{\text{sin-cos}}(\{\hat{\chi_i}\}_{i=1}^{N_{\chi}}, \{\chi_i\}_{i=1}^{N_{\chi}}) =
\frac{1}{2 N_\chi} \sum_{i=1}^{N_{\chi}} (\cos \hat{\chi_i}  - \cos \chi_i )^2 +(\sin \hat{\chi_i}
- \sin \chi_i)^2
\label{sincosLoss}
\EE
Notably, this loss function is justified by a nice geometric interpretation, whereby it is
equivalent to computing the cosine loss between the 2D vectors that describe the $\chi$ angles on
the unit circle (proof in Eq.~\ref{eqn:proof_of_equivalence_sin_cos}).


{\bf (iii) Normal vectors to the dihedral plane.} $\chi$ angles are examples of dihedral angles,
meaning that they are defined as the angle between two planes. For $\chi$ angles, the two planes
are described by subsequent triplets of atoms along the side-chains. Any two subsequent $\chi$
angles share one plane. Therefore, any conformation with $N_{\chi}$ angles can be alternatively
described by $N_{\chi}+1$ planes (or their normal vectors); one of these normal vectors is a
redundant internal coordinate (defined by backbone + C$\beta$ atoms), while others specify the
$N_\chi$ independent degrees of freedom.

We consider training models to predict the dihedral planes' normal vectors:
$\vn_{\chi_1}\,...\,\vn_{\chi_4}$. It should be noted that unlike the sine/cosine transformation,
the vectors are not invariant to rotations, but \textit{equivariant} of type $\ell = 1$ (geometric
vectors) which  can be extracted from the H-Packer equivariant network. We use a cosine loss over
the true and predicted vectors:
\EQ
\mathcal{L}_{\text{norms}}(\{\hat{\vn}_{\chi_i}\}_{i=1}^{N_{\chi}},
\{\vn_{\chi_i}\}_{i=1}^{N_{\chi}}) = \frac{1}{N_\chi} \sum_{i=1}^{N_{\chi}} 1 - \langle
\hat{\vn}_{\chi_i}, \vn_{\chi_i} \rangle
\label{planeloss}
\EE

{\bf Relevant symmetries in computing loss functions.}
Some amino acid conformations exhibit a rotation symmetry by $\pi$ in some of their $\chi$ angles.
For example, $\chi_2$ of Phenylalanine and Tyrosine indicates the torsion of their benzene rings,
thus a rotation by $\pi$ leaves the conformation physically unchanged. However, as $\chi$ angles
are formally defined by internal atom names, these equivalent conformations are associated with
different $\chi$ angle values. We correct for this degeneracy by considering the minimum loss value
between considering $\chi$ and $\pi - \chi$ as targets during training and evaluation. When
computing the error on the atomic coordinates (generally via Root Mean Square Deviation, RMSD) for
the full side-chain, we need to consider other such symmetries between non-$\chi$ atoms, as listed
in Table~\ref{table:symmetric_conformations}.


\section{Results}

\begin{figure}
    \centering
    \resizebox{\textwidth}{!}{
    \begin{tabular}{c c}
        \includegraphics[valign=m,width=0.85\textwidth]{Quals/figures/simple_task/simple_task_fixed_symmetries__test_predictions_by_model_type.pdf}
        &
        \includegraphics[valign=m,width=0.15\textwidth]{Quals/figures/simple_task/lmax_legend.pdf}
    \end{tabular}
    }
    \figcaption{Test MAE for the simple task of predicting $\chi$ angles from atomic
    conformation.}{Panels show reconstruction accuracies using three loss functions: the angle $\chi$ itself
    (left), the sin/cos transform of the angle (center), and the normal vectors to the dihedral
    planes (right), for different maximum angular degrees $\ell_\text{max}$ (colors).}
    \label{fig:simple_task_performance_by_model_type}
    \vspace{-4mm} % reduces whitespace after the table by desired amount!!!
\end{figure}

\subsection{Toy task: inferring $\chi$ angles from atomic coordinates}
\label{sec:simple_task}
We start by studying the behavior of our model on a simple task: predicting (or rather,
calculating) $\chi$ angles from the true atomic coordinates of the conformation. We found this to
be a useful benchmark to study our model's behavior.

\textbf{Setup.} We randomly select 160 structures from our real task's training set (see below) and
split them into 100/30/30 for training/validation/testing, respectively. We collect conformations
of all residues presenting $\chi$ angles, and consider only their heavy atoms (C, N, O, S). We then
apply the Zernike encoding varying $\ell_{\text{max}}$ from $1$ to $5$ and train models with
varying $\ell_{\text{max}}$ consistent with that of the input, as well as with different prediction
objectives (angles, sin-cos of angles, plane norms). Crucially, we vary the number of hidden
channels (decreasing it with higher $\ell_{\text{max}}$) to keep the number of parameters constant
around 330k, and thus, removing differences in model capacity as a contributing factor to
performance. We do not condition the models with amino-acid identity to make the problem more
challenging, and therefore more interesting. We refer to Section~\ref{sec:training_details} for
more details.

\textbf{Results.} Test Mean Absolute Error (MAE) per $\chi$ angle for all models is shown in
Figure~\ref{fig:simple_task_performance_by_model_type}, and training curves are shown in the
Appendix (Figure~\ref{fig:simple_task_training_curves}). Notably, the Angle model performs the
worst, and is unable to recover the true $\chi$ angle with negligible error. The Sin-Cos and Plane
Norm models instead recover all $\chi$ angles with very low error (< 5 \AA) with $\ell_{\text{max}}
> 1$. 
 It appears that $\ell_{\text{max}}=2$ is the minimum sufficient degree nedded to solve this task
 with high accuracy. We note that error is higher for later $\chi$ angles. We hypothesise that this
 is expected for two reasons: (i) later $\chi$ angles depend on atoms that are farther way from the
 center of the neighborhood, thus having lower angular resolution within the Zernike
 representation, and (ii) there is simply less training data for them.  Weighting $\chi$ angles in
 the loss function according to their average frequency partially mitigates the second issue
 (Figure~\ref{fig:simple_task_performance_by_model_type_with_weighted_loss}). 
 Notably, the fact that the model performs well without explicit knowledge of amino acid identities
 implies that it can easily infer the amino acid type from the  the number and the relative
 location of the atoms.

% \textbf{Takeaways.} Based on these results, we 

\begin{table}[t]
    \centering
    \resizebox{0.9\textwidth}{!}{
    \begin{tabular}{l c c c c | c c c | c c c}
     \multicolumn{1}{l|}{\bf CASP13} & \multicolumn{4}{c}{Angle MAE $\degree$ $\downarrow$} &
     \multicolumn{3}{c}{Angle Accuracy \% $\uparrow$} & \multicolumn{3}{c}{Atom RMSD \AA
     $\downarrow$} \\
     % \hline
     \cline{1-1}
     \textbf{Method} & $\chi_1$ & $\chi_2$ & $\chi_3$ & $\chi_4$ & All & Core & Surface & All &
     Core & Surface \\
     \hline
     \hline
     SCWRL         & 27.64 & 28.97 & 49.75 & 61.54  &  56.2 & 71.3 & 43.4  &  0.934 & 0.495 & 1.027
     \\
     FASPR         & 27.04 & 28.41 & 50.30 & 60.89  &  56.4 & 70.3 & 43.6  &  0.910 & 0.502 & 1.002
     \\
     RosettaPacker & 25.88 & 28.25 & 48.13 & 59.82  &  58.6 & 75.3 & 35.7  &  0.872 & 0.422 & 1.001
     \\
     DLPacker      & 22.18 & 27.00 & 51.22 & 70.04  &  58.8 & \underline{73.9} & 45.4  &  0.772 &
     0.402 & 0.876 \\
     AttnPacker    & \underline{18.92} & \underline{23.17} & \underline{44.89} & 58.98  & 
     \underline{62.1} & 73.7 & \underline{47.6}  &  \underline{0.669} & \underline{0.366} &
     \underline{0.775} \\
     DiffPack      & \textbf{15.35} & \textbf{19.19} & \textbf{37.30} & \textbf{50.19}  & 
     \textbf{69.5} & \textbf{82.7} & \textbf{57.3}  &  \textbf{0.579} & \textbf{0.298} &
     \textbf{0.696} \\
     \hline
     \hline
     % H-Packer$_0$ $\ell_{\text{max}} = 4$ & 27.60 & 32.31 & 47.87 & 52.77 & 48.3 & 59.1 & 40.5 & 0.979 & 0.771 & 1.136 \\
     % H-Packer$_2$ $\ell_{\text{max}} = 4$ & 24.36 & 29.94 & 46.01 & 52.89 & 53.0 & 67.7 & 42.5 & 0.884 & 0.610 & 1.083 \\
     % H-Packer$_5$ $\ell_{\text{max}} = 4$ & 24.27 & 29.76 & 46.32 & 52.56 & 53.5 & 68.8 & 43.0 & 0.878 & 0.594 & 1.081 \\
     % H-Packer$_{up}$ $\ell_{\text{max}} = 4$ & \textit{20.78} & \textit{27.48} & \textit{44.58} & \textit{51.97} & \textit{56.8} & \textit{72.7} & \textit{45.6} & \textit{0.790} & \textit{0.514} & \textit{0.999} \\
     H-Packer$_0^{\ell_{\text{max}}=5}$    & 26.89 & 31.95 & 47.51 & 52.75  &  49.3 & 61.5 & 40.4 
     &  0.961 & 0.726 & 1.131 \\[1.4pt]
     H-Packer$_2^{\ell_{\text{max}}=5}$     & 23.64 & 29.47 & 45.17 & 53.26  &  54.4 & 69.9 & 43.6 
     &  0.863 & 0.575 & 1.070 \\[1.4pt]
     H-Packer$_5^{\ell_{\text{max}}=5}$     & 23.60 & 29.40 & \underline{44.91} & \underline{52.91}
      &  54.7 & 70.7 & 43.7  &  0.858 & 0.564 & 1.067 \\[1.4pt]
     \hline
     H-Packer$_{up}^{\ell_{\text{max}}=5}$  & \textit{20.03} & \textit{26.88} & \textit{42.74} &
     \textit{52.07}  &  \textit{58.4} & \textit{75.4} & \textit{46.4}  &  \textit{0.765} &
     \textit{0.483 }& \textit{0.980} \\
     \hline
    \end{tabular}
    }
    \vspace{6mm}

    \resizebox{0.9\textwidth}{!}{
    \begin{tabular}{l c c c c | c c c | c c c}
     \multicolumn{1}{l|}{\bf CASP14} & \multicolumn{4}{c}{Angle MAE $\degree$ $\downarrow$} &
     \multicolumn{3}{c}{Angle Accuracy \% $\uparrow$} & \multicolumn{3}{c}{Atom RMSD \AA
     $\downarrow$} \\
     % \hline
     \cline{1-1}
     \textbf{Method} & $\chi_1$ & $\chi_2$ & $\chi_3$ & $\chi_4$ & All & Core & Surface & All &
     Core & Surface \\
     \hline
     \hline
     SCWRL         & 33.50 & 33.05 & 51.61 & 55.28  &  45.4 & 62.5 & 33.2  &  1.062 & 0.567 & 1.216
     \\
     FASPR         & 33.04 & 32.49 & 50.15 & 54.82  &  46.3 & 62.4 & 34.0  &  1.048 & 0.594 & 1.205
     \\
     RosettaPacker & 31.79 & 28.25 & 50.54 & 56.16  &  47.5 & \underline{67.2} & 33.5  &  1.006 &
     0.501 & 1.183 \\
     DLPacker      & 29.01 & 33.00 & 53.98 & 72.88  &  48.0 & 66.9 & 33.9  &  0.929 & 0.476 & 1.107
     \\
     AttnPacker    & \underline{25.34} & \underline{28.19} & 48.77 & \underline{51.92}  & 
     \underline{50.9} & 66.2 & 36.3  &  \underline{0.823} & \underline{0.438} & \underline{1.001} \\
     DiffPack      & \textbf{21.91} & \textbf{25.54} & \textbf{44.27} & 55.03  &  \textbf{57.5} &
     \textbf{77.8} & \textbf{43.5}  &  \textbf{0.770} & \textbf{0.356} & \textbf{0.956} \\
     \hline
     % \hline
     % H-Packer$_0$ $\ell_{\text{max}} = 4$ & 33.14 & 36.10 & 49.30 & 50.22 & 39.8 & 55.9 & 30.7 & 1.102 & 0.780 & 1.307 \\
     % H-Packer$_2$ $\ell_{\text{max}} = 4$ & 30.43 & 34.49 & 48.13 & 50.39 & 43.3 & 63.0 & 32.7 & 1.027 & 0.658 & 1.265 \\
     % H-Packer$_5$ $\ell_{\text{max}} = 4$ & 30.36 & 34.38 & 48.76 & 50.62 & 43.5 & 64.1 & 32.5 & 1.024 & 0.648 & 1.260 \\
     % H-Packer$_{up}$ $\ell_{\text{max}} = 4$ & \textit{26.84} & \textit{32.10} & \textit{46.90} & \textit{50.03} & \textit{46.7} & \textit{68.1} & \textit{35.0} & \textit{0.934} & \textit{0.557} & \textit{1.178} \\
     % \hline
     H-Packer$_0^{\ell_{\text{max}}=5}$    & 32.31 & 35.90 & 49.05 & 50.34 & 40.8 & 57.0 & 31.6 &
     1.087 & 0.762 & 1.297 \\[1.4pt]
     H-Packer$_2^{\ell_{\text{max}}=5}$    & 29.96 & 34.32 & 47.46 & 50.50 & 45.0 & 65.1 & 34.1 &
     1.011 & 0.629 & 1.250 \\[1.4pt]
     H-Packer$_5^{\ell_{\text{max}}=5}$    & 29.61 & 34.03 & \underline{46.72} & \textbf{50.35} &
     45.2 & 65.5 & 34.0 & 1.002 & 0.626 & 1.244 \\[1.4pt]
     \hline
     H-Packer$_{up}^{\ell_{\text{max}}=5}$ & \textit{26.58} & \textit{31.54} & \textit{45.67} &
     \textit{49.46} & \textit{48.1} & \textit{69.5} & \textit{36.2} & \textit{0.915} &
     \textit{0.534 }& \textit{1.160} \\
     \hline
    \end{tabular}
    }
    \figcaption{Comparative assessment of H-Packer on CASP13 and CASP14.}{
    We present best results in \textbf{bold} and second-best \underline{underlined}.
    \textit{Italicized} results represent an upper bound to our algorithm's performance.
    Performance of models other than H-Packer is taken from~\citep{zhang_diffpack_2023}.}
    \label{tab:results_casp13_and_14}
\vspace{-6mm} % reduces whitespace after the table by desired amount!!!
\end{table}


\subsection{Side-Chain Packing}

\textbf{Dataset.} We consider the training and validation datasets used in 
DLpacker~\citep{misiura_dlpacker_2022}, consisting of 19,436 structures with a maximum
inter-protein sequence similarity of 50\%. Unlike DLPacker, we do not remodel structures with
PDB-redo~\citep{joosten_pdb_redo_2009} and do not convert selenomethionine residues into
methionine. For testing our model, we use the CASP13 and CASP14 targets (82 and 64 structures,
respectively). We remove from the training and validation sets any protein that has sequence
similarity above 50\% with any of the proteins in the test set.

\textbf{H-Packer training.} We used the Sin/Cos loss function (Eq.~\ref{sincosLoss}) as it was the
best-performing loss in our toy-task; while the Plane Norms loss (Eq.~\ref{planeloss}) also
performed well in the toy-task, we found that models trained with the Sin/Cos objective were easier
to regularize via dropout in the final invariant feed-forward neural network. The {\em initial
guess} and the {\em refinement} networks were trained with the same $\ell_{\text{max}}$ of 5 and
$n_{\text{max}}=12$; the latter was chosen such that it included at least one radial function with
wavelength lower than the minimum interatomic distance. We also considered models trained with
$\ell_{\text{max}} = 4$, tuning the number of hidden features to keep the number of trainable
parameters the same as the $\ell_{\text{max}} = 5$ models, and equal to $\sim$3M. All models were
trained for 10 epochs, keeping the model with lowest validation loss at the end of an epoch; see
further details in Section~\ref{sec:training_details}. %In preliminary experiments, we briefly
experimented with up-weighting higher-index $\chi$ angles in the loss function
(Section~\ref{sec:simple_task}), but ultimately decided against it in an effort to promote more
accurate learning of the lower-index ones $\chi_1$ and $\chi_2$.
Throughout our experiments, we consider the performance of H-Packer models with different number of
rounds of refinement. For example, H-Packer$_0$ denotes the model with {\em no} refinement. For
each model, we also compute an {\em upper bound} in performance of the refinement process by
tasking the \textit{refinement} model to predict $\chi$ angles from the ground truth neighboring
structures (i.e., the toy task); we denote this by H-Packer$_{up}$.

\textbf{Metrics.} In line with previous work~\citep{mcpartlon_end--end_2023, zhang_diffpack_2023},
we evaluate our models on three main metrics. (i) Angle-specific Mean Absolute Error (MAE), (ii)
residue-level angle accuracy, defined as the proportion of residues for which the prediction of all
$\chi$ angles is within $20\degree$ of the true value, and (iii) average atomic Root Mean Square
Deviation (RMSD) of side-chain atoms across residues. We further distinguish between
\textit{Surface} and  \textit{Core} residues, as conformations occurring on the surface of proteins
are notoriously harder to predict. Surface residues are defined as having at most 15 $\beta$-C
within 10~\AA\, of their $\beta$-C, whereas core residues must have at least 20 $\beta$-C's in this
range.


\begin{table}[t]
    \begin{center}
    \resizebox{0.9\textwidth}{!}{
    \begin{tabular}{l c c c c | c c c c}
     % & \multicolumn{8}{c}{Atom RMSD (All) \AA $\downarrow$} \\
     & \multicolumn{4}{c}{CASP13} & \multicolumn{4}{c}{CASP14} \\
     % \hline
     \textbf{Method} & base & Std. & Sym. & Std.+Sym. & base & Std. & Sym. & Std.+Sym. \\
     \hline
     \hline
     % H-Packer$_0$ $\ell_{\text{max}} = 4$    & 0.979 & 0.946 & 0.943 & 0.927  &  1.102 & 1.086 & 1.065 & 1.049 \\
     % H-Packer$_2$ $\ell_{\text{max}} = 4$    & 0.884 & 0.863 & 0.848 & 0.828  &  1.027 & 1.010 & 0.989 & 0.973 \\
     % H-Packer$_5$ $\ell_{\text{max}} = 4$    & 0.878 & 0.858 & 0.841 & 0.821  &  1.024 & 1.007 & 0.986 & 0.969 \\
     % H-Packer$_{up}$ $\ell_{\text{max}} = 4$ & 0.790 & 0.767 & 0.755 & 0.733  &  0.934 & 0.915 & 0.899 & 0.881 \\
     % \hline
     H-Packer$_0^{\ell_{\text{max}} = 5}$    & 0.961 & 0.943 & 0.923 & 0.906  &  1.087 & 1.070 &
     1.050 & 1.034 \\[1.4pt]
     H-Packer$_2^{\ell_{\text{max}} = 5}$    & 0.863 & 0.842 & 0.826 & 0.805  &  1.011 & 0.992 &
     0.972 & 0.953 \\[1.4pt]
     H-Packer$_5^{\ell_{\text{max}} = 5}$    & 0.858 & 0.837 & 0.821 & 0.800  &  1.002 & 0.984 &
     0.964 & 0.945 \\[1.4pt]
     \hline
     H-Packer$_{up}^{\ell_{\text{max}} = 5}$ & 0.765 & 0.741 & 0.730 & 0.706  &  0.915 & 0.895 &
     0.880 & 0.860 \\
     \hline
    \end{tabular}
    }
    \end{center}
    \figcaption{Atom RMSD (\AA $\downarrow$) for H-Packer across all residues with different
    treatments of the true structure.}{\textbf{Std:} standardizing the true structure by reconstructing
    it with our data-derived redundant internal coordinates and the real $\chi$ angles. \textbf{Sym:} 
    considering the additional non-natural symmetries used by AttnPacker.}
    \label{tab:rmsd_inflation}
    \vspace{2mm} % adds/reduces whitespace after the table by desired amount!!!
\end{table}

\begin{table}[t]
    \centering
    \resizebox{1.0\textwidth}{!}{
    \begin{tabular}{l l c c c c | c c c | c c c}
     & & \multicolumn{4}{c}{Angle MAE $\degree$ $\downarrow$} & \multicolumn{3}{c}{Angle Accuracy
     \% $\uparrow$} & \multicolumn{3}{c}{Atom RMSD \AA $\downarrow$} \\
     % \hline
     & \textbf{H-Packer$_5$} & $\chi_1$ & $\chi_2$ & $\chi_3$ & $\chi_4$ & All & Core & Surface &
     All & Core & Surface \\
     \hline
     \hline
     \multirow{2}{*}{CASP13}
     & $\ell_{\text{max}} = 4$ & 24.27 & 29.76 & 46.32 & \textbf{52.56} & 53.5 & 68.8 & 43.0 &
     0.878 & 0.594 & 1.081 \\
     & $\ell_{\text{max}} = 5$ & \textbf{23.60} & \textbf{29.40} & \textbf{44.91} & 52.91 &
     \textbf{54.7} & \textbf{70.7} & \textbf{43.7} & \textbf{0.858} & \textbf{0.564} &
     \textbf{1.067} \\
     \hline
     \multirow{2}{*}{CASP14}
     & $\ell_{\text{max}} = 4$ & 30.36 & 34.38 & 48.76 & 50.62 & 43.5 & 64.1 & 32.5 & 1.024 & 0.648
     & 1.260 \\
     & $\ell_{\text{max}} = 5$ & \textbf{29.61} & \textbf{34.03} & \textbf{46.72} & \textbf{50.35}
     & \textbf{45.2} & \textbf{65.5} & \textbf{34.0} & \textbf{1.002} & \textbf{0.626} &
     \textbf{1.244} \\
     \hline
    \end{tabular}
    }
    \figcaption{Ablation in $\ell_{\text{max}}$ of H-Packer.}{Metrics for the other H-packer models 
    can be found in Table~\ref{tab:full_ablation_in_lmax}.}
    \label{tab:lmax_ablation}
% \vspace{-4mm} % reduces whitespace after the table by desired amount!!!
\end{table}


\textbf{Comparative Evaluation on CASP13 and CASP14 targets.}
Table~\ref{tab:results_casp13_and_14} compare H-Packer's performance in side-chain packing with
other computational methods~\citep{mcpartlon_end--end_2023, zhang_diffpack_2023,
misiura_dlpacker_2022, krivov_improved_2009, alford_rosetta_2017, huang_faspr_2020}. Despite its
simplicity, H-Packer$_5$ is competitive against the state-of-the-art at predicting $\chi_3$ and
$\chi_4$, but falls behind on $\chi_1$ and $\chi_2$ predictions. This discrepancy indicates that
H-Packer has likely learned complementary features to the other models. Moreover, H-Packer mostly
outperforms the physics-based computational algorithms~\citep{krivov_improved_2009,
alford_rosetta_2017, huang_faspr_2020} and is competitive with
DLPacker~\citep{misiura_dlpacker_2022} in all our performance metrics. Interestingly, H-Packer is
consistently better than physics-based approaches in terms of overall Atom RMSD, but tends to fall
shorter on Angle Accuracy. We present error distrubtions for H-Packer in
Figures~\ref{fig:casp13_error_distributions_part_1},~\ref{fig:casp13_error_distributions_part_2},~\ref{fig:casp14_error_distributions_part_1}, and~\ref{fig:casp14_error_distributions_part_2}.
 
Interestingly, while H-Packer predictions are improved upon using refinement networks, the
performance saturates after 2 steps of refinement; the accuracies after 5 iterations of refinement
are comparable to those after only 2 steps (Table~\ref{tab:results_casp13_and_14}). Therefore, it
is unlikely that further refinement could improve H-Packer's performance to reach its upper bound
performance. We hypothesise that training H-Packer to produce confidence scores might help in
developing site-specific convergence criteria to help bridge the
gap~\citep{mcpartlon_end--end_2023, zhang_diffpack_2023}.

Our experiments show that H-packer is competitive against physics-based methods and some
machine-learning solutions, but its performance still lags behind the state-of-the-art at
predicting $\chi$ angles closer to the backbone. Overall, the lack of consistent comparative
patterns in performance metrics suggests that H-Packer learns features complementary to other
approaches.

\textbf{Ablation in $\ell_{\text{max}}$.} Table~\ref{tab:lmax_ablation} shows how changing
$\ell_{\text{max}}$ (from 4 to 5) impacts  the performance of H-Packer. For the same model
capacity, using higher $\ell_{\text{max}}$ consistently yields better performance, indicating that
higher angular resolutions of the input can be beneficial for learning this task. This performance
improvement comes with a trade-off in  training and inference time, which scale superlinearly with
$\ell_{\text{max}}$ unless the Tensor Product computation is adequately
constrained~\citep{cobb_efficient_2021}. For reference, training  our models with
$\ell_{\text{max}} = 5$ takes $\sim$40\% longer than those with $\ell_{\text{max}} = 4$. We leave
the hyperparameter optimization of   $\ell_{\text{max}}$ to future work. 

\textbf{On computing RMSD fairly.} In Table~\ref{tab:results_casp13_and_14} we report RMSD computed
by measuring the distance between the coordinates of true and predicted atoms, modulo the
symmetries we report in~\ref{table:symmetric_conformations}. However, other algorithms such as
AttnPacker~\citep{mcpartlon_end--end_2023} consider other symmetries as well, sometimes even
between atoms of differing chemical elements. Though these symmetries reflect spatially similar
conformations (such as a flip of the Histidine ring), they result in inflated RMSD scores. We show
the effect of this inflation on H-Packer predictions in Table~\ref{tab:rmsd_inflation}. In the same
table, we also show the RMSD computed against true structures that have been "reconstructed" using
the true $\chi$ angles and the constant values that we use for redundant internal coordinate within
H-Packer; we do this in an effort to disentangle the Null Reconstruction Error
(Figure~\ref{fig:null_reconstruction}) from the error given by mistakes in $\chi$ angle prediction.

\textbf{Speed.} Table~\ref{tab:speed} shows relative reconstruction speeds for several packing
algorithms. Using the current implementation of the reconstruction algorithm, the best-performing
H-Packer model is about 7x faster than the popular algorithm RosettaPacker and 6x faster than
DLPacker; however, it is considerably slower than AttnPacker. Speed can be considerably cut down by
half at the expense of minor performance degradation using two refinement iterations instead of
five. However, more considerable speed gains may be achieved by CPU parallelization, when computing
 holographic encodings of structural neighborhood during initial data processing. Indeed, each
initial guessing and refinement step of H-Packer predicts all $\chi$ angles at once, but in the
current implementation holographic encodings are computed in series, creating a bottleneck that
currently accounts for ~88\% of the inference time (10\% is atom placement, and only 2\% is making
the actual predictions on GPU). We plan on optimizing this aspect in future iterations of the model.

\begin{table}[t]
    \centering
    \resizebox{1.0\textwidth}{!}{
    \begin{tabular}{l c c c c c c c c}
        \textbf{Method} & HPacker$_{5}^{\ell_{\text{max}}=5}$ & HPacker$_{2}^{\ell_{\text{max}}=5}$
        & HPacker$_{0}^{\ell_{\text{max}}=5}$ & AttnPack & DLPack & RosPack & FASPR & SCWRL4 \\
        \hline
        \textbf{Rel. Time} & 1.00 & 0.51 & 0.18 & 0.05 & 5.70 & 6.96 & 0.02 & 0.67 \\
        \hline
    \end{tabular}
    }
    \figcaption{Relative times to undertake full atomic reconstruction.}{In our current (unoptimized)
    implementation, HPacker$_{5}^{\ell_{\text{max}}=5}$ takes 1,482s
    to reconstruct the 82 CASP13 targets on a single NVIDIA A40 GPU. Times for the other methods
    were taken from ~\citep{mcpartlon_end--end_2023}.}
    \label{tab:speed}
\vspace{-6mm} % reduces whitespace after the table by desired amount!!!
\end{table}



